%% SC-STEM: design notes for the simulation study.
%% Internal document, standalone, NOT included in main.tex or supplement.tex.
%% Section 10 describes two development extensions that are neither in the
%% package nor in the paper, and must not be mentioned there.
\documentclass[11pt]{article}
\usepackage[margin=2.5cm]{geometry}
\usepackage{amsmath,amssymb,amsthm}
\usepackage{booktabs}
\usepackage{array}
\usepackage{enumitem}
\usepackage[round]{natbib}
\usepackage[hidelinks]{hyperref}

\newcommand{\E}{\mathbb{E}}
\newcommand{\Var}{\mathrm{Var}}
\newcommand{\Cov}{\mathrm{Cov}}
\newcommand{\Corr}{\mathrm{Corr}}
\newcommand{\KL}{\mathrm{KL}}
\newcommand{\Stem}{\textsf{Stem}}
\newtheorem{proposition}{Proposition}
\newtheorem{remark}{Remark}

\title{The SC-STEM simulation study:\\ design, generator and separation of the regimes}
\author{Design notes, internal}
\date{27 September 2026}

\begin{document}
\maketitle

\begin{abstract}
\noindent These notes fix the design of the simulation study of the SC-STEM paper: what the study has to show, the generator shared by every scenario, the two variants of dependence between regimes, the baseline regime, the scenarios and their parameter values, and the relation of each scenario to the data-generating processes of SC-STEM and of STEM. They then derive a measure of how separated the regimes are as seen by the assignment step of the algorithm, first formally and then numerically for every scenario. Section~\ref{sec:dev} records two development extensions, constrained estimation with blocks of parameters common across regimes and a coupled model with dependence between regimes; they are implemented outside the package, in the development area of the repository, and are not mentioned in the paper.
\end{abstract}

\tableofcontents

%%=============================================================================
\section{What the study has to show}\label{sec:aims}
%%=============================================================================

SC-STEM admits \emph{completely heterogeneous} regimes: each regime is a spatio-temporal model with its own full set of parameters. The study starts from there and asks three questions.

\begin{enumerate}[leftmargin=2em]
  \item \textbf{Complete heterogeneity.} When the data are generated by SC-STEM with every block of parameters different across regimes, does the procedure identify the correct model and estimate it well? ``Identify'' has three targets: the number of regimes $k$, the partition of the locations (which location belongs to which regime), and the regime-specific parameters $\Psi_g$. The question is asked at two degrees of separation of the regimes, poorly and well separated.
  \item \textbf{Complete homogeneity.} When the data are generated by the pooled STEM model, with a single regime, how often does the procedure declare more regimes than there are? This is the false-positive rate.
  \item \textbf{A common component.} When one block of parameters (the regression coefficients, the persistence, the innovation variance, the spatial range, or the whole measurement error) is common to all regimes and the others differ, what changes?
\end{enumerate}

Two remarks qualify the first and the third question.

\begin{remark}[Which blocks are heterogeneous]
The package estimates every block regime by regime. It therefore has no likelihood-based way of deciding \emph{which} blocks differ across regimes; the only instrument it offers is the bootstrap test of the between-regime differences. A likelihood-based answer requires the constrained estimator of Section~\ref{sec:dev:constr}, which is a development extension.
\end{remark}

\begin{remark}[What a common block can and cannot do for the free estimator]\label{rem:common}
Given the partition, the package fits each regime on its own data only. The sampling distribution of $\hat\Psi_g$ then depends only on the data and the parameters of regime $g$, and not on whether a block of $\Psi_g$ happens to equal the corresponding block of $\Psi_h$. A block that is common in the data-generating process therefore does \emph{not} make the other parameters better estimated by the free estimator: its only effect is to remove one source of contrast between the regimes, which can only make the partition harder to recover. The gain in precision the third question is after exists only for an estimator that \emph{knows} the block is common, that is the constrained estimator of Section~\ref{sec:dev:constr}. Since that estimator is not part of the paper, the third question is read in the paper as \emph{how much each block contributes to the separation of the regimes}; the comparison between the free and the constrained estimator is kept in the development area.
\end{remark}

%%=============================================================================
\section{The model and its special case}\label{sec:model}
%%=============================================================================

\subsection{SC-STEM}

Let $g_i \in \{1,\dots,k\}$ be the regime of location $i$ and $I_g = \{i : g_i = g\}$. For $i \in I_g$ and $t = 1,\dots,T$,
\begin{align}
  z_{it} &= x_{it}^\top \beta_g + K_i\, y^{(g)}_t + e_{it}, \label{eq:meas}\\
  y^{(g)}_t &= G_g\, y^{(g)}_{t-1} + \eta^{(g)}_t, \qquad \eta^{(g)}_t \sim N(0, \Sigma_{\eta,g}), \label{eq:state}\\
  e_{t,I_g} &\sim N\bigl(0, \Sigma_{e,g}\bigr), \qquad \Sigma_{e,g} = \sigma^2_{\epsilon,g} I + \sigma^2_{\omega,g}\, C_{\theta_g}, \quad [C_{\theta}]_{ij} = e^{-\theta h_{ij}}, \label{eq:err}
\end{align}
with $e_t$ independent over $t$, and with the latent processes and the error fields of different regimes independent of one another. Each regime carries its own
\[
  \Psi_g = (\beta_g,\ G_g,\ \Sigma_{\eta,g},\ \sigma^2_{\epsilon,g},\ \sigma^2_{\omega,g},\ \theta_g,\ m_{0,g}),
\]
and, because the regimes are independent, the log-likelihood of a partition is the sum of the regime log-likelihoods,
\begin{equation}\label{eq:loglik}
  \ell(\Psi_1,\dots,\Psi_k) = \sum_{g=1}^k \ell_g(\Psi_g),
\end{equation}
each $\ell_g$ being the STEM log-likelihood of the locations in $I_g$ alone. SC-STEM \emph{allows} every component of $\Psi_g$ to differ across regimes; it does not require them to differ, and a configuration in which some components coincide is a point of its parameter space like any other.

\subsection{STEM as the special case with one regime}

With $k = 1$ there is one latent process, one error field over the whole network and one set of parameters: \eqref{eq:meas}--\eqref{eq:err} reduce to the pooled STEM model. STEM is SC-STEM at $k = 1$.

\subsection{Two sources of separation}

Under \eqref{eq:meas}--\eqref{eq:err} the regimes differ in two ways. They differ in their parameters $\Psi_g$; and they differ because each has its own latent path $y^{(g)}$ and its own error field, independent of those of the other regimes. The second is a property of the model, not of the parameters, and it is written in the covariance of the data: conditionally on the covariates,
\begin{equation}\label{eq:blockcov}
  \Cov(z_{it}, z_{js}) = 0 \qquad \text{for } i \in I_g,\ j \in I_h,\ g \neq h,
\end{equation}
so that $\Cov(z_t)$ is block-diagonal by regime. Two locations in different regimes are uncorrelated however close they are. Both sources recur throughout: the first is what the scenarios vary, the second is what the variants of Section~\ref{sec:variants} switch on and off.

\subsection{What each step of the algorithm sees}\label{sec:score}

The alternating algorithm has two steps, and they do not see the same things.

\paragraph{Parameter step.} Given the partition, each regime is fitted by the EM algorithm on the exact likelihood $\ell_g$. It sees every component of $\Psi_g$, the spatial correlation included.

\paragraph{Assignment step.} Given the parameters, location $i$ is scored against regime $h$ by the pseudo-likelihood
\begin{equation}\label{eq:score}
  S_i(h) = \sum_{t=1}^T \left[ -\tfrac12 \log\bigl(2\pi s^2_h\bigr) - \frac{\bigl(z_{it} - x_{it}^\top\beta_h - K_i\hat y^{(h)}_t\bigr)^2}{2 s^2_h} \right],
  \qquad s^2_h = \sigma^2_{\epsilon,h} + \sigma^2_{\omega,h},
\end{equation}
where $\hat y^{(h)}_t$ is the smoothed latent path of regime $h$, and the label maximizes $S_i(h)$ plus the Potts term. The score sees $\beta_h$, the total error variance $s^2_h$, and the dynamics $(G_h, \Sigma_{\eta,h})$ only through the smoothed path $\hat y^{(h)}$. It does \emph{not} see the range $\theta_h$, nor the split of $s^2_h$ between nugget and partial sill, nor any spatial correlation of the errors, within or across regimes. Those enter only through the exact likelihood of the parameter step. A contrast in $\theta$ therefore separates the regimes much less, for the assignment step, than a contrast of the same nominal size in $\beta$.

%%=============================================================================
\section{The common generator}\label{sec:generator}
%%=============================================================================

Every scenario is generated by one generator. A scenario is a choice of the parameter sets $\Psi_1, \Psi_2, \Psi_3$; a variant is a choice of the dependence between regimes (Section~\ref{sec:variants}). The number of regimes is $K = 3$ in every scenario except the pooled one, where the regimes are replaced by a single one.

\subsection{Locations}

The locations lie in the plane and distances are Euclidean in its units. The regime sizes $n_1, n_2, n_3$ are equal up to the remainder (balanced) or proportional to $1:2:3$ (unbalanced), with no regime below six locations. The coordinates follow the overlap design,
\begin{equation}\label{eq:coords}
  s_i \mid g_i = g \ \sim\ N_2\bigl(\mu_g,\ \nu_{sp}\, I_2\bigr), \qquad \nu_{sp} = \nu_{\mathrm{tot}} - \Var(\mu),
\end{equation}
with the three centres $\mu_g$ at the vertices of an equilateral triangle of side $2\omega$ and $\nu_{\mathrm{tot}} = 0.4 + 2/3$ the total variance of a coordinate, held fixed so that the network covers the same area, about four units across, whatever $\omega$ is. The overlap $\omega$ controls how mixed the regimes are in space: at $\omega = 0$ the three clouds coincide and the partition has no spatial signature; at the reference value $\omega = 0.70$ the centres are $1.63$ within-regime standard deviations apart, so the regimes are contiguous areas that interpenetrate along their borders. In the pooled scenario there are no regimes, and the locations form one cloud, $s_i \sim N_2(0, \nu_{\mathrm{tot}} I_2)$, covering the same area: the network is then not spatially clustered, so the selection rule is given no geometric cue for telling one regime from several.

\subsection{Covariate}

One exogenous covariate, an AR(1) in time whose innovations are a spatially correlated field:
\begin{equation}\label{eq:cov}
  x_t = a\, x_{t-1} + \sqrt{1-a^2}\; u_t, \qquad u_t \sim N_n(0, C_x), \qquad [C_x]_{ij} = e^{-3 h_{ij}/R_x},
\end{equation}
with $a = 0.7$ (a shock to the covariate halves in about two periods) and practical range $R_x = 4$, the whole network, so the covariate varies on a large spatial scale. The covariate is then centred and scaled over all $nT$ values, so that $\E x \approx 0$ and $\Var x \approx 1$. Its spatial correlation sits between the two degenerate cases: a covariate independent across locations would make a coefficient contrast trivially visible, one common to all locations would be indistinguishable from the latent process. The design matrix is $x_{it} = (1, x_{it})^\top$ and $\beta_g = (\beta_{0,g}, \beta_{1,g})^\top$.

\subsection{Latent processes}

One AR(1) per regime, with loading $K_i = 1$:
\begin{equation}\label{eq:latent}
  y^{(g)}_t = G_g\, y^{(g)}_{t-1} + \eta^{(g)}_t, \qquad
  \eta^{(g)}_t = \sigma_{\eta,g}\left(\sqrt{\rho}\,\xi_t + \sqrt{1-\rho}\,\zeta_{g,t}\right),
\end{equation}
with $\xi_t$ and $\zeta_{g,t}$ independent standard normal, independent over $t$ and $g$. Then $\Var(\eta^{(g)}_t) = \sigma^2_{\eta,g}$ and $\Corr(\eta^{(g)}_t, \eta^{(h)}_t) = \rho$ for $g \neq h$: $\xi_t$ is the shock common to all regimes and $\zeta_{g,t}$ the one specific to regime $g$. The processes start from zero and run through a burn-in of 200 periods that is discarded, which puts $y_1$ at the stationary distribution to within $\max_g G_g^{200}$.

\paragraph{Stationary moments.} Writing $y^{(g)}_t = \sum_{j \ge 0} G_g^j \eta^{(g)}_{t-j}$,
\begin{align}
  v_g &:= \Var\bigl(y^{(g)}_t\bigr) = \frac{\sigma^2_{\eta,g}}{1 - G_g^2}, \label{eq:vg}\\
  \Cov\bigl(y^{(g)}_t, y^{(h)}_t\bigr) &= \sum_{j\ge0} (G_gG_h)^j\, \rho\,\sigma_{\eta,g}\sigma_{\eta,h}
     = \frac{\rho\,\sigma_{\eta,g}\sigma_{\eta,h}}{1 - G_gG_h}, \label{eq:cross}\\
  \Corr\bigl(y^{(g)}_t, y^{(h)}_t\bigr) &= \rho\,\kappa_{gh}, \qquad
     \kappa_{gh} = \frac{\sqrt{(1-G_g^2)(1-G_h^2)}}{1 - G_gG_h} \le 1, \label{eq:kappa}
\end{align}
where the sums use the independence of the innovations over time. Since $(1 - G_gG_h)^2 - (1-G_g^2)(1-G_h^2) = (G_g - G_h)^2 \ge 0$, $\kappa_{gh} \le 1$, with equality if and only if $G_g = G_h$: two autoregressions with different persistence cannot be perfectly correlated, even when they are driven by the same shocks. The dynamics of a regime are described by the pair $(G_g, v_g)$, the persistence and the amplitude of the latent signal, and $\sigma^2_{\eta,g} = v_g(1 - G_g^2)$ follows.

\subsection{Measurement error}

Two constructions, independent over $t$.
\begin{description}[leftmargin=1.5em]
  \item[Fields by regime.] Each regime draws its own field, independently of the others:
  \begin{equation}\label{eq:blockfield}
    e_{t,I_g} \sim N\bigl(0,\ \sigma^2_{\epsilon,g} I + \sigma^2_{\omega,g} C_{\theta_g}\bigr), \qquad g = 1,2,3.
  \end{equation}
  Two locations in different regimes have uncorrelated errors, however close they are.
  \item[One field.] One spatial field spans the whole network, scaled by regime:
  \begin{equation}\label{eq:onefield}
    e_{it} = \sigma_{\epsilon,g_i}\, \varepsilon_{it} + \sigma_{\omega,g_i}\, w_t(s_i), \qquad
    \varepsilon_{it} \sim N(0,1)\ \text{i.i.d.}, \qquad w_t \sim N_n(0, C_\theta).
  \end{equation}
  Then $\Cov(e_{it}, e_{jt}) = \sigma_{\omega,g_i}\sigma_{\omega,g_j}\, e^{-\theta h_{ij}}$ also across a regime border. A single field has a single range, so this construction requires $\theta$ to be common to all regimes; the nugget and the partial sill may still differ by regime.
\end{description}
Within each regime the two constructions give the same marginal law, $e_{t,I_g} \sim N(0, \sigma^2_{\epsilon,g} I + \sigma^2_{\omega,g} C_\theta)$: they differ only in the covariance \emph{between} regimes.

\subsection{Response}

\begin{equation}\label{eq:resp}
  z_{it} = \beta_{0,g_i} + \beta_{1,g_i}\, x_{it} + y^{(g_i)}_t + e_{it}.
\end{equation}
The conditional mean given the latent path, $\mu_{it} = \beta_{0,g_i} + \beta_{1,g_i} x_{it} + y^{(g_i)}_t$, is recorded as the target of predictive accuracy.

\subsection{Parametrization}

Each regime is described by
\[
  \bigl(\beta_{0,g},\ \beta_{1,g},\ G_g,\ v_g,\ \sigma^2_{\epsilon,g},\ \sigma^2_{\omega,g},\ R_g\bigr),
  \qquad \sigma^2_{\eta,g} = v_g(1-G_g^2), \qquad \theta_g = 3/R_g,
\]
where $R_g$ is the practical range, the distance at which the spatial correlation of the error falls to $e^{-3} \approx 5\%$. The five blocks of parameters used throughout are
\[
  \text{beta} = (\beta_0, \beta_1), \quad G, \quad \Sigma_\eta = \sigma^2_\eta, \quad \theta, \quad \text{error} = (\sigma^2_\epsilon, \sigma^2_\omega, \theta).
\]

%%=============================================================================
\section{Dependence between regimes: independent and shared}\label{sec:variants}
%%=============================================================================

\subsection{Two separate assumptions}

The generator has two links between regimes, and they are \emph{separate assumptions on different components of the model}: neither implies the other.
\begin{itemize}[leftmargin=1.5em]
  \item $\rho$ governs the latent processes: whether the regimes share their shocks.
  \item The error construction governs the measurement error: whether one field crosses the regime borders.
\end{itemize}
They are paired by a choice of design, so that each scenario has an ``all independent'' and an ``all shared'' variant, and the two mixed combinations are added where they exist.

\subsection{The variants}

\begin{center}
\begin{tabular}{ll>{\raggedright\arraybackslash}p{5.4cm}>{\raggedright\arraybackslash}p{4.6cm}}
\toprule
code & $\rho$ & error & exists when \\
\midrule
\texttt{ind} & 0 & fields by regime & always (except the pooled scenario) \\
\texttt{shr} & 1 & one field if $\theta$ is common, otherwise fields by regime & always \\
\texttt{lat} & 1 & fields by regime & only when $\theta$ is common (otherwise it is \texttt{shr}) \\
\texttt{fld} & 0 & one field & only when $\theta$ is common \\
\bottomrule
\end{tabular}
\end{center}

\paragraph{Independent} (\texttt{ind}) is exactly the SC-STEM model: independent latent paths, block-diagonal error covariance, the likelihood \eqref{eq:loglik}.

\paragraph{Shared} (\texttt{shr}) shares everything the scenario permits. With $\rho = 1$ the regimes are driven by one sequence of shocks $\xi_t$, which each regime filters with its own persistence $G_g$ and scales by its own $\sigma_{\eta,g}$; the correlation of the latent paths is then $\kappa_{gh}$ of \eqref{eq:kappa}. The content of ``shared'' therefore depends on the scenario: where the ranges differ (complete heterogeneity, and the scenarios with a common block other than $\theta$ or the error) only the latent can be shared, since one field has one range; where $\theta$ is common, the latent and the field are both shared. In the pooled scenario everything is shared by definition, and there is only one variant.

\paragraph{The mixed combinations} (\texttt{lat}, \texttt{fld}) separate the two links where both can be switched, which is only in the scenarios with a common $\theta$. Elsewhere \texttt{lat} coincides with \texttt{shr}, and \texttt{fld} does not exist.

\subsection{What SC-STEM makes of the shared variants}

Under \texttt{shr}, \texttt{lat} and \texttt{fld} the marginal law of each regime taken alone is still exactly its STEM model: a Gaussian vector restricted to a subset of its components is Gaussian with the corresponding sub-covariance, so the error of regime $g$ has covariance $\sigma^2_{\epsilon,g} I + \sigma^2_{\omega,g} C_\theta$ on $I_g$; and each $y^{(g)}$ is marginally an AR(1) with parameters $(G_g, \sigma^2_{\eta,g})$. SC-STEM is wrong only in its assumption that the regimes are independent. Four consequences follow.
\begin{enumerate}[leftmargin=2em]
  \item \textbf{Estimation.} The sum \eqref{eq:loglik} becomes a composite likelihood, the product of the block marginals of the true joint law. Given the partition, each $\ell_g$ is the exact likelihood of the data of regime $g$, so $\hat\Psi_g$ remains consistent.
  \item \textbf{Joint uncertainty.} The parametric bootstrap of the package regenerates the regimes independently, and so ignores the correlation between $\hat\Psi_g$ and $\hat\Psi_h$ that the shared variants induce. This concerns the tests of the between-regime differences.
  \item \textbf{The partition.} The independence of the regimes is itself information about the partition (Section~\ref{sec:sep}). Sharing the latent weakens the channel through $\hat y^{(g)}$ that the assignment score \eqref{eq:score} uses; sharing the field changes nothing in the score, which does not see spatial correlation, and acts only through the exact likelihood of the parameter step and the comparison across $k$.
  \item \textbf{The choice of $k$.} The criteria compare $\sum_g \ell_g$ across $k$, and that sum ignores the dependence between the blocks.
\end{enumerate}
The shared variants therefore test the robustness of SC-STEM to a mild violation of its assumptions, one that leaves each regime correctly specified.

\subsection{Empirical readings}

Both variants describe real situations, and they differ in what the regime border is.
\begin{description}[leftmargin=1.5em]
  \item[Independent: the border is a barrier.] The regimes are distinct systems with drivers of their own, and the correlation stops at the border. Examples: river basins separated by a watershed divide; hydraulically disconnected aquifers; epidemic outbreaks in areas with no mobility between them. This is the case the model is written for and the most favourable to it, but in real data it is the less common one.
  \item[Shared: the border is a change of response, not a barrier.] All regimes are subject to the same large-scale driver, and each filters it with its own persistence and amplitude, while local disturbances spread across the border. Examples: temperature anomalies on a coast and inland, driven by the same synoptic weather but with different persistence; regional unemployment responding to the same national cycle; soil moisture, where the rain is common and the soil type decides how long a shock lasts; a land-use boundary, across which the error field is continuous. This is the more common case in real data.
\end{description}

\subsection{The pooled likelihood against the likelihood of a partition}\label{sec:mi}

When the data come from the pooled STEM model, the SC-STEM log-likelihood of any partition into $k > 1$ regimes is the sum of block marginals, while the pooled log-likelihood is the joint. At the true parameters,
\begin{equation}\label{eq:mi}
  \E\Bigl[\log p(z) - \sum_{g} \log p(z_{I_g})\Bigr] = \mathcal{I}\bigl(z_{I_1}; \dots; z_{I_k}\bigr) \ \ge\ 0,
\end{equation}
the multi-information (total correlation) between the blocks, which is zero only if the blocks are independent. Splitting a pooled network into regimes therefore \emph{loses} likelihood, besides paying the penalty of the criterion, and the loss grows with $T$ and with the strength of the dependence the blocks share: the variance of the common latent process and the range of the common field. False positives are then rare when the common dependence is strong; the case at risk is a weak common dependence, where splitting costs almost nothing and the additional parameters fit noise. The false-positive rate measured in the pooled scenario is thus a function of its latent variance and range, which are those of the baseline regime.

%%=============================================================================
\section{The baseline regime}\label{sec:base}
%%=============================================================================

Regime 1 of every scenario, and the only regime of the pooled scenario, carries the baseline:
\begin{center}
\begin{tabular}{ll>{\raggedright\arraybackslash}p{9.2cm}}
\toprule
parameter & value & reading \\
\midrule
$\beta_0$ & 2 & mean level, two standard deviations above zero (coefficient of variation 0.5) \\
$\beta_1$ & $0.55 \approx \sqrt{0.30}$ & the covariate explains 30\% of the variance \\
$G$ & 0.8 & a latent shock halves in $\log 0.5/\log 0.8 = 3.1$ periods \\
$v$ & 0.30 & the latent process explains 30\% of the variance; $\sigma^2_\eta = 0.108$ \\
$\sigma^2_\epsilon$ & 0.20 & nugget, 20\% of the variance \\
$\sigma^2_\omega$ & 0.20 & partial sill, 20\% of the variance \\
$R$ & 2 & practical range, half the network; $\theta = 1.5$ \\
\bottomrule
\end{tabular}
\end{center}
The variance of $z$ at a location of regime 1 is $\beta_1^2 + v + \sigma^2_\epsilon + \sigma^2_\omega = 0.30 + 0.30 + 0.20 + 0.20 \approx 1$, so every contrast below reads in standard deviations of the response. The systematic part (covariate and latent) explains 60\%, the local part (field and nugget) 40\%, and the total error variance is $s^2 = 0.40$.

%%=============================================================================
\section{The scenarios}\label{sec:scen}
%%=============================================================================

\subsection{S1: complete heterogeneity}

Every block differs across regimes. Regime 1 is the baseline, regime 3 carries the full contrast, and regime 2 sits halfway.
\begin{center}
\footnotesize
\setlength{\tabcolsep}{4pt}
\begin{tabular}{l ccc ccc >{\raggedright\arraybackslash}p{3.9cm}}
\toprule
& \multicolumn{3}{c}{S1-weak (poorly separated)} & \multicolumn{3}{c}{S1-strong (well separated)} & \\
\cmidrule(lr){2-4}\cmidrule(lr){5-7}
parameter & $g=1$ & $g=2$ & $g=3$ & $g=1$ & $g=2$ & $g=3$ & reading \\
\midrule
$\beta_0$ & 2.00 & 2.10 & 2.20 & 2.00 & 2.30 & 2.60 & level step of 0.1 or 0.3 sd \\
$\beta_1$ & 0.55 & 0.65 & 0.75 & 0.55 & 0.85 & 1.15 & effect $+18\%$ per step, or doubled \\
$G$ & 0.80 & 0.70 & 0.60 & 0.80 & 0.50 & 0.20 & half-life 3.1/1.9/1.4, or 3.1/1.0/0.4 \\
$v$ & 0.30 & 0.35 & 0.40 & 0.30 & 0.40 & 0.50 & the latent weighs more \\
$\sigma^2_\eta$ & 0.108 & 0.179 & 0.256 & 0.108 & 0.300 & 0.480 & $= v(1-G^2)$ \\
$\sigma^2_\epsilon$ & 0.20 & 0.25 & 0.30 & 0.20 & 0.35 & 0.50 & more unstructured noise \\
$\sigma^2_\omega$ & 0.20 & 0.175 & 0.15 & 0.20 & 0.15 & 0.10 & less spatial noise \\
$R$ & 2 & 1.5 & 1 & 2 & 1 & 0.5 & range from half the network to an eighth \\
\midrule
$s^2$ & 0.40 & 0.425 & 0.45 & 0.40 & 0.50 & 0.60 & total error variance \\
nugget share & 0.50 & 0.59 & 0.67 & 0.50 & 0.70 & 0.83 & $\sigma^2_\epsilon/s^2$ \\
\bottomrule
\end{tabular}
\end{center}
In words: along the regimes the level rises, the covariate matters more, shocks die out faster but the latent signal is larger, and the error becomes noisier, less spatially structured and more local. In the strong level the variance of $z$ in regime 3 is $1.15^2 + 0.50 + 0.60 = 2.42$, against 1 in regime 1. The variants are \texttt{ind} and \texttt{shr}; since the ranges differ, \texttt{shr} shares the latent only.

\emph{Question:} does the procedure recover $k = 3$, the partition and $\Psi_1,\Psi_2,\Psi_3$ of the model written in full?

\subsection{S2: complete homogeneity}

The pooled STEM model at the baseline:
\begin{equation}\label{eq:pooled}
  z_{it} = 2 + 0.55\, x_{it} + y_t + e_{it}, \qquad y_t = 0.8\, y_{t-1} + \eta_t,\ \ \sigma^2_\eta = 0.108, \qquad
  e_t \sim N\bigl(0,\ 0.20\, I + 0.20\, C_{1.5}\bigr),
\end{equation}
one latent path and one field over the whole network. In the generator it is the corner in which the three parameter sets coincide, $\rho = 1$ (so the three latent paths coincide, since equal dynamics and one sequence of shocks give one path) and the field is one. There is one variant only: with $\rho = 0$ or fields by regime it would no longer be the pooled model.

\emph{Question:} how often does the procedure declare $k > 1$? By \eqref{eq:mi}, the answer depends on the strength of the common dependence, which here is that of the baseline.

\subsection{S3: one common block}

One block takes the baseline value in all regimes; the other blocks keep the values of S1-strong, so that the partition is recoverable and the comparison with S1-strong isolates the block.
\begin{center}
\begin{tabular}{l>{\raggedright\arraybackslash}p{3.6cm}>{\raggedright\arraybackslash}p{8.4cm}}
\toprule
scenario & common & by regime \\
\midrule
S3-beta  & $\beta = (2, 0.55)$ & everything else as S1-strong \\
S3-G     & $G = 0.8$ & $v$ as S1-strong ($0.30/0.40/0.50$), hence $\sigma^2_\eta = 0.108/0.144/0.180$ \\
S3-Seta  & $\sigma^2_\eta = 0.108$ & $G$ as S1-strong ($0.8/0.5/0.2$), hence $v = 0.300/0.144/0.113$ \\
S3-theta & $R = 2$ ($\theta = 1.5$) & $\sigma^2_\epsilon$, $\sigma^2_\omega$ and everything else as S1-strong \\
S3-error & $\sigma^2_\epsilon = \sigma^2_\omega = 0.20$, $R = 2$ & $\beta$, $G$, $v$ as S1-strong \\
\bottomrule
\end{tabular}
\end{center}
The rule for the dynamics follows the parametrization by $(G, v)$: when $G$ is common the amplitudes $v_g$ stay those of S1-strong; when $\sigma^2_\eta$ is common the persistences $G_g$ stay those of S1-strong and the amplitudes follow from \eqref{eq:vg}.

S3-beta, S3-G and S3-Seta have the variants \texttt{ind} and \texttt{shr} (latent only). S3-theta and S3-error have a common range and so all four: \texttt{ind}, \texttt{shr} (latent and field), \texttt{lat} and \texttt{fld}.

\emph{Question, as read in the paper:} how much does each block contribute to the separation of the regimes? By Remark~\ref{rem:common} the free estimator does not become more precise when a block is common; the question of precision belongs to the constrained estimator of Section~\ref{sec:dev:constr}.

\subsection{The scenario-variants}

\begin{center}
\begin{tabular}{llcl}
\toprule
scenario & variants & count & shared components in \texttt{shr} \\
\midrule
S1-weak, S1-strong & \texttt{ind}, \texttt{shr} & $2 \times 2$ & latent \\
S2 & one & 1 & everything \\
S3-beta, S3-G, S3-Seta & \texttt{ind}, \texttt{shr} & $3 \times 2$ & latent \\
S3-theta, S3-error & \texttt{ind}, \texttt{shr}, \texttt{lat}, \texttt{fld} & $2 \times 4$ & latent and field \\
\midrule
total & & 19 & \\
\bottomrule
\end{tabular}
\end{center}
The mixed combinations cost four scenario-variants out of nineteen. Each scenario-variant is then crossed with the levels of the other factors of the design (Section~\ref{sec:design}) and replicated.

\subsection{Relation to the data-generating processes of SC-STEM and STEM}

\begin{center}
\begin{tabular}{ll>{\raggedright\arraybackslash}p{5.6cm}>{\raggedright\arraybackslash}p{3.9cm}}
\toprule
scenario & variant & SC-STEM, \eqref{eq:meas}--\eqref{eq:err} & STEM \\
\midrule
\textbf{S1} & \texttt{ind} & \textbf{exact: the reference case} & no: one regime cannot represent it \\
S1 & \texttt{shr} & no: each regime is exact, the latents are correlated across regimes & no \\
\textbf{S2} & (one) & exact at $k = 1$ & \textbf{exact: the special case} \\
S3-x & \texttt{ind} & exact & no \\
S3-x & \texttt{shr}, \texttt{lat}, \texttt{fld} & no: each regime is exact, the regimes are dependent (latent, field or both) & no \\
\bottomrule
\end{tabular}
\end{center}

\paragraph{The reference case} is S1 in the variant \texttt{ind}: the model \eqref{eq:meas}--\eqref{eq:err} written in full, every block different and the regimes independent.

\paragraph{The interpretable special case} is S2: the pooled STEM, which is SC-STEM at $k = 1$.

\paragraph{Why S3-x in the variant \texttt{ind} is exact.} SC-STEM lets the parameters differ across regimes but does not require them to; a configuration in which one block coincides is a point of its parameter space. The free estimator is correctly specified there, and estimates the common block regime by regime, spending more parameters than necessary.

\paragraph{Making every block common does not give the pooled model.} If every block of $\Psi_g$ coincides but the regimes remain independent, the data are three independent copies of one STEM model on three sub-networks: three latent paths and three fields, and the covariance is still block-diagonal \eqref{eq:blockcov}. The pooled model is \emph{common parameters, plus one latent path ($\rho = 1$), plus one field}. The chain of S3 therefore does not end in S2, and the partition of three independent copies remains identifiable from the dependence alone (Section~\ref{sec:sep}).

%%=============================================================================
\section{Separation of the regimes: theory}\label{sec:sep}
%%=============================================================================

The parameter values of Section~\ref{sec:scen} are meant to produce poorly and well separated regimes. This section defines what ``separated'' means for the algorithm and derives a measure of it; Section~\ref{sec:sepnum} evaluates it.

\subsection{The quantity}

Take a location $i$ of regime $g$ and a wrong regime $h \neq g$. Ignoring the Potts term, the assignment step prefers $g$ to $h$ when $S_i(g) - S_i(h) > 0$, with $S_i$ the score \eqref{eq:score}. The measure of separation is the expected value of that difference,
\begin{equation}\label{eq:gap}
  \E\bigl[S_i(g) - S_i(h)\bigr] = \sum_{t=1}^T \E\bigl[d_{it}\bigr], \qquad
  d_{it} = \log \phi\bigl(z_{it};\, m^{(g)}_{it},\, s^2_g\bigr) - \log \phi\bigl(z_{it};\, m^{(h)}_{it},\, s^2_h\bigr),
\end{equation}
where $\phi(\cdot\,; m, s^2)$ is the normal density and $m^{(r)}_{it} = x_{it}^\top\beta_r + \hat y^{(r)}_t$ is the mean the score attributes to location $i$ under regime $r$.

\subsection{Assumptions}

\begin{enumerate}[label=(A\arabic*), leftmargin=2.5em]
  \item \emph{Oracle paths and parameters.} The parameters of every regime are at their true values and the smoothed paths equal the true ones, $\hat y^{(r)}_t = y^{(r)}_t$. This is the situation the algorithm reaches when the regimes are well estimated; it is an idealization, and the measure is therefore the separation available to the assignment step at best.
  \item \emph{Covariate.} $x_{it}$ is independent of the latent processes and of the errors, with $\E x_{it} = 0$ and $\E x_{it}^2 = 1$ (the covariate is standardized over all $nT$ values, so these hold on average over the network rather than exactly at each location).
  \item \emph{Stationarity.} The latent processes are at their stationary distribution, with the moments \eqref{eq:vg}--\eqref{eq:cross}.
  \item \emph{Errors.} $e_{it}$ has mean zero and variance $s^2_g$, and is independent of $x$ and of the latent processes.
\end{enumerate}

\subsection{Derivation}

Under the true regime, the residual is the error itself: $z_{it} - m^{(g)}_{it} = e_{it}$ by (A1) and \eqref{eq:resp}. Under the wrong regime,
\begin{equation}\label{eq:resid}
  z_{it} - m^{(h)}_{it} = \underbrace{x_{it}^\top(\beta_g - \beta_h)}_{a_{it}} + \underbrace{y^{(g)}_t - y^{(h)}_t}_{b_t} + e_{it}.
\end{equation}
Writing out the two log densities,
\[
  d_{it} = -\tfrac12 \log s^2_g - \frac{e_{it}^2}{2 s^2_g} + \tfrac12 \log s^2_h + \frac{(a_{it} + b_t + e_{it})^2}{2 s^2_h}.
\]
Taking expectations, $\E e_{it}^2 = s^2_g$ by (A4), and
\[
  \E(a_{it} + b_t + e_{it})^2 = \E a_{it}^2 + \E b_t^2 + \E e_{it}^2,
\]
because the three cross products have zero expectation: $\E[a_{it} e_{it}] = 0$ and $\E[b_t e_{it}] = 0$ by (A4), and $\E[a_{it} b_t] = \E[a_{it}]\,\E[b_t] = 0$ by (A2) and $\E b_t = 0$. Hence
\begin{equation}\label{eq:D}
  D_{gh} := \E\bigl[d_{it}\bigr] = \frac12\left[\frac{s^2_g + \Delta^2_{gh}}{s^2_h} - 1 - \log\frac{s^2_g}{s^2_h}\right],
  \qquad \Delta^2_{gh} = \E a_{it}^2 + \E b_t^2 .
\end{equation}
The two components of $\Delta^2_{gh}$ follow from (A2) and (A3). With $x_{it} = (1, x_{it})$,
\begin{equation}\label{eq:mean}
  \E a_{it}^2 = \E\bigl(\Delta\beta_0 + \Delta\beta_1 x_{it}\bigr)^2 = \Delta\beta_0^2 + 2\Delta\beta_0\Delta\beta_1\E x_{it} + \Delta\beta_1^2 \E x_{it}^2 = \Delta\beta_0^2 + \Delta\beta_1^2,
\end{equation}
with $\Delta\beta_j = \beta_{j,g} - \beta_{j,h}$; and, from \eqref{eq:vg}--\eqref{eq:kappa},
\begin{equation}\label{eq:lat}
  \E b_t^2 = v_g + v_h - 2\,\Cov\bigl(y^{(g)}_t, y^{(h)}_t\bigr) = v_g + v_h - \frac{2\rho\,\sigma_{\eta,g}\sigma_{\eta,h}}{1 - G_gG_h}
  = v_g + v_h - 2\rho\,\kappa_{gh}\sqrt{v_g v_h}.
\end{equation}
Equations \eqref{eq:D}, \eqref{eq:mean} and \eqref{eq:lat} are the three formulas of the measure.

\paragraph{Reading.} $D_{gh}$ is the expected Kullback--Leibler divergence between the two normal densities the score compares at one time point,
\[
  D_{gh} = \E_{x,y}\Bigl[\KL\bigl(N(m^{(g)}_{it}, s^2_g)\ \big\|\ N(m^{(h)}_{it}, s^2_h)\bigr)\Bigr],
\]
since $\KL(N(\mu_1,\sigma_1^2)\|N(\mu_2,\sigma_2^2)) = \tfrac12[\sigma_1^2/\sigma_2^2 + (\mu_1-\mu_2)^2/\sigma_2^2 - 1 - \log(\sigma_1^2/\sigma_2^2)]$ and $\E(m^{(g)}_{it} - m^{(h)}_{it})^2 = \Delta^2_{gh}$. It is non-negative, zero if and only if $s^2_g = s^2_h$ and $\Delta^2_{gh} = 0$, and not symmetric in $g$ and $h$. It separates into a mean part, $\Delta^2_{gh}/(2s^2_h)$, and a variance part, $\tfrac12[s^2_g/s^2_h - 1 - \log(s^2_g/s^2_h)]$, which is the only place the error variances enter.

\subsection{Why one time point, and why multiplying by $T$}

The score \eqref{eq:score} is a sum over time of one term per period: it treats the periods as conditionally independent given the latent path, as the STEM error is white in time. $D_{gh}$ is the expectation of one of those terms. By linearity of expectation, and because under (A2)--(A4) every period has the same expectation,
\begin{equation}\label{eq:T}
  \E\bigl[S_i(g) - S_i(h)\bigr] = \sum_{t=1}^T \E[d_{it}] = T\, D_{gh}.
\end{equation}
This step uses only linearity: it does not require the terms to be independent over time, and they are not, since $x$ and the latent paths are autocorrelated. Independence would matter for the variance of the difference, not for its mean. $T\, D_{gh}$ is in units of log-likelihood (nats): it is the expected log-likelihood ratio in favour of the true regime that one location accumulates over its whole series, and $e^{T D_{gh}}$ the corresponding likelihood ratio. It grows linearly in $T$, while its standard deviation grows more slowly, so a location is assigned more reliably the longer its series. For the algorithm the relevant pair is the hardest one: with three regimes placed on a line of contrasts, the adjacent pairs $(1,2)$ and $(2,3)$ in both directions, and the minimum over them.

\subsection{Formal consequences}

\begin{proposition}[Independence alone separates the regimes]\label{prop:indep}
Under $\rho = 0$, $\E b_t^2 = v_g + v_h$ whatever the parameters, so
\[
  D_{gh} \ \ge\ \frac12\left[\frac{s^2_g + v_g + v_h}{s^2_h} - 1 - \log\frac{s^2_g}{s^2_h}\right] \ \ge\ \frac{v_g + v_h}{2 s^2_h} \quad \text{when } s^2_g = s^2_h .
\]
\end{proposition}
With independent regimes the latent paths separate the regimes on their own, with a bound set by the weight of the latent signal relative to the noise and not by any contrast in the parameters. In the variant \texttt{ind}, therefore, the level of separation of the parameters (weak or strong) governs the estimation of $\Psi_g$ and the choice of $k$, but hardly the recovery of the partition, which the latent channel already secures. This is coherent with \texttt{ind} being the reference case: the model is correct and the problem is easy. The only lever that makes the independent case hard for the partition is the latent weight $v$ relative to $s^2$, which belongs to the baseline.

\begin{proposition}[Decomposition of the latent term]
\[
  \E b_t^2 = \bigl(\sqrt{v_g} - \sqrt{v_h}\bigr)^2 + 2\sqrt{v_g v_h}\,\bigl(1 - \rho\,\kappa_{gh}\bigr).
\]
\end{proposition}
The first term is the difference in amplitude, the second the difference in timing, which is due to independent shocks ($\rho < 1$) or to different persistence ($\kappa_{gh} < 1$). With $\rho = 1$ and a common $G$ the second term vanishes and only the amplitude separates the paths: the paths are perfectly correlated and differ only in scale. With $\rho = 1$, a common $G$ and a common $v$ the latent term vanishes altogether, and the separation rests on $\beta$ and $s^2$ alone.

\begin{proposition}[What the measure does not see]
$D_{gh}$ depends on the error only through the total variances $s^2_g, s^2_h$. It does not depend on $\theta$, on the split of $s^2$ into nugget and partial sill, nor on whether the error is drawn by regime or as one field.
\end{proposition}
Two consequences for the design. The mixed variants have the same measure as their counterpart with the same $\rho$: \texttt{lat} as \texttt{shr}, \texttt{fld} as \texttt{ind}. And S3-theta has the same measure as S1-strong, since the two differ only in $\theta$. Whatever the field variant or the range contribute to recovery, they contribute through the exact likelihood of the parameter step, not through the assignment step.

\subsection{Limits of the measure}

The measure describes the assignment step at the oracle: in practice $\hat y^{(h)}$ is estimated from the locations currently in $h$, which early in the alternation include misassigned ones, and the parameters are estimates. It ignores the Potts term, which adds to the score of a label the agreement with the neighbours. It is pairwise and per location. It measures what the assignment step can use, not the total information in the data, which the parameter step reads through the exact likelihood. It is a guide to calibrate the scenarios, not a prediction of the recovery rates.

%%=============================================================================
\section{Separation of the regimes: numbers}\label{sec:sepnum}
%%=============================================================================

Table~\ref{tab:sep} evaluates \eqref{eq:D}--\eqref{eq:lat} at the parameter values of Section~\ref{sec:scen}, for the four ordered adjacent pairs, at $\rho = 0$ (variants \texttt{ind} and \texttt{fld}) and $\rho = 1$ (variants \texttt{shr} and \texttt{lat}). The last columns give, for the hardest pair, the mean term $\Delta\beta_0^2 + \Delta\beta_1^2$, the latent term $\E b_t^2$, the total $T\,D$ at $T = 120$, and the correlation $\rho\kappa$ of the latent paths of the adjacent regimes.

\begin{table}[ht]
\centering
\small
\setlength{\tabcolsep}{4.5pt}
\caption{The separation measure $D_{gh}$ per time point for the adjacent pairs, and its components at the hardest pair. S2 has $D = 0$ by construction.}\label{tab:sep}
\begin{tabular}{lc cccc c cc c cc}
\toprule
scenario & $\rho$ & $D_{12}$ & $D_{21}$ & $D_{23}$ & $D_{32}$ & $\min D$ & mean & latent & $120\,\min D$ & $\rho\kappa_{12}$ & $\rho\kappa_{23}$ \\
\midrule
S1-weak   & 0 & 0.789 & 0.838 & 0.856 & 0.907 & 0.789 & 0.02 & 0.650 & 94.7  & 0     & 0     \\
          & 1 & 0.047 & 0.050 & 0.037 & 0.040 & 0.037 & 0.02 & 0.013 & 4.5   & 0.974 & 0.985 \\
S1-strong & 0 & 0.892 & 1.113 & 0.908 & 1.089 & 0.892 & 0.18 & 0.700 & 107.0 & 0     & 0     \\
          & 1 & 0.292 & 0.363 & 0.205 & 0.246 & 0.205 & 0.18 & 0.057 & 24.6  & 0.866 & 0.943 \\
S3-beta   & 0 & 0.712 & 0.888 & 0.758 & 0.909 & 0.712 & 0    & 0.700 & 85.4  & 0     & 0     \\
          & 1 & 0.112 & 0.138 & 0.055 & 0.066 & 0.055 & 0    & 0.057 & 6.6   & 0.866 & 0.943 \\
S3-G      & 0 & 0.892 & 1.113 & 0.908 & 1.089 & 0.892 & 0.18 & 0.700 & 107.0 & 0     & 0     \\
          & 1 & 0.199 & 0.247 & 0.163 & 0.194 & 0.163 & 0.18 & 0.006 & 19.5  & 1     & 1     \\
S3-Seta   & 0 & 0.636 & 0.793 & 0.372 & 0.445 & 0.372 & 0.18 & 0.257 & 44.6  & 0     & 0     \\
          & 1 & 0.276 & 0.343 & 0.172 & 0.205 & 0.172 & 0.18 & 0.017 & 20.6  & 0.866 & 0.943 \\
S3-theta  & 0 & 0.892 & 1.113 & 0.908 & 1.089 & 0.892 & 0.18 & 0.700 & 107.0 & 0     & 0     \\
          & 1 & 0.292 & 0.363 & 0.205 & 0.246 & 0.205 & 0.18 & 0.057 & 24.6  & 0.866 & 0.943 \\
S3-error  & 0 & 1.100 & 1.100 & 1.350 & 1.350 & 1.100 & 0.18 & 0.700 & 132.0 & 0     & 0     \\
          & 1 & 0.350 & 0.350 & 0.296 & 0.296 & 0.296 & 0.18 & 0.057 & 35.5  & 0.866 & 0.943 \\
\bottomrule
\end{tabular}
\end{table}

\paragraph{How to read a number.} Take S1-strong at $\rho = 1$: the hardest pair is regime 2 scored against regime 3, with $D_{23} = 0.205$. The mean term is $0.3^2 + 0.3^2 = 0.18$; the latent term is $0.40 + 0.50 - 2 \times 0.943 \sqrt{0.40 \times 0.50} = 0.057$, small because the two paths, driven by the same shocks, are correlated at 0.94; the variances are $s^2_2 = 0.50$ and $s^2_3 = 0.60$. Then $D_{23} = \tfrac12[(0.50 + 0.237)/0.60 - 1 - \log(0.50/0.60)] = 0.205$ per period, and a location of regime 2 accumulates on average $120 \times 0.205 = 24.6$ nats in favour of its own regime over a series of 120 periods, a likelihood ratio of about $e^{24.6}$.

\paragraph{The independent variant: the latent separates on its own.} At $\rho = 0$ the latent term alone is $v_g + v_h$, between $0.26$ and $0.70$, larger than any mean term of the design, and the measure at $T = 120$ lies between 45 and 132 nats in every scenario. This is Proposition~\ref{prop:indep} in numbers: the difference between S1-weak (94.7) and S1-strong (107.0) is small, because the parameters add little to what the independence of the regimes already provides. The lowest value, S3-Seta (44.6), is the one where a common innovation variance with fast dynamics shrinks the latent amplitudes to $v = 0.144$ and $0.113$, weakening the latent channel.

\paragraph{The shared variant: the two levels work as intended.} At $\rho = 1$ the latent term almost vanishes (0.013 in S1-weak, 0.057 in S1-strong) and the separation rests on the parameters. S1-weak is genuinely poorly separated, 4.5 nats per location over the series, where the Potts term and the neighbours can decide the label; S1-strong is well separated, 24.6 nats.

\paragraph{S3 in the shared variant: what each block contributes.} Against S1-strong (24.6):
\begin{itemize}[leftmargin=1.5em]
  \item \emph{S3-beta} falls to 6.6: the regression coefficients are the main separator for the assignment step. What is left comes from the variances and from the residual timing difference of the latent paths.
  \item \emph{S3-G} falls to 19.5: with a common persistence and common shocks the paths are perfectly correlated ($\rho\kappa = 1$) and differ only in amplitude, $(\sqrt{0.40} - \sqrt{0.50})^2 = 0.006$.
  \item \emph{S3-Seta} falls to 20.6: different persistence keeps the paths apart in timing, but their amplitudes are smaller.
  \item \emph{S3-theta} stays at 24.6, identical to S1-strong: the score does not see the range. Whatever S3-theta loses or gains in recovery against S1-strong is due to the parameter step alone.
  \item \emph{S3-error} rises to 35.5, above S1-strong: the common total variance is the baseline $0.40$, the smallest of S1-strong, so every mean contrast is measured against a smaller noise, and the variance part, which was the only place the error entered, drops to zero.
\end{itemize}

%%=============================================================================
\section{Design of the experiment}\label{sec:design}
%%=============================================================================

\paragraph{Factors.} The scenario-variants are crossed with five margins: the number of locations $n$, the length of the series $T$, the overlap $\omega$, the balance of the regime sizes and the number of neighbours $m$ of the graph of the Potts penalty. The reference levels are $n = 100$, $T = 120$, $\omega = 0.70$, balanced regimes and $m = 5$. The design is organized in blocks: each block is the full factorial of the margins it lists, every other margin at its reference level, crossed with the scenario-variants, and the design is the union of the blocks, a cell that two blocks share being run once. Blocks and levels are lists in the setup of the script. In S2 the overlap and the balance are vacuous and the scenario enters once per $(n, T, m)$. No network smaller than $n = 40$ is used: the grid reaches $k = 4$ regimes of at least six locations each, which needs $n \ge 24$, and $n = 40$ already gives about 13 locations per regime.

\paragraph{Scenario-variants for the paper, and additional ones.} Nine scenario-variants answer the questions the paper asks: S2 (false positives), S1-strong-\texttt{ind} (the reference case), S1-strong-\texttt{shr} and S1-weak-\texttt{shr} (the two levels of separation with coupled regimes), and the five S3-x-\texttt{shr} (the contribution of each block, in the variant where the separation rests on the parameters). The other ten are additional: S1-weak-\texttt{ind} and the five S3-x-\texttt{ind}, which at the reference sit at the ceiling because the latent paths separate the regimes on their own, and the four \texttt{lat} and \texttt{fld} variants, which take the mechanism apart channel by channel.

\paragraph{The experiment.} The nine scenario-variants of the paper in four blocks. The core crosses the three factors that govern recovery; the other three vary one robustness margin at a time around the reference: the balance, the graph, and a network at the national scale.
\begin{center}
\begin{tabular}{l>{\raggedright\arraybackslash}p{5.4cm}>{\raggedright\arraybackslash}p{4.6cm}r}
\toprule
block & margins (others at the reference) & cells & \\
\midrule
core & $n \in \{40, 100, 200\}$, $T \in \{60, 120, 365\}$, $\omega \in \{0, 0.70, 1\}$ & $8 \times 27 + 9$ (S2 per $(n, T)$) & 225 \\
balance & unbalanced ($1:2:3$) & 8 (S2 vacuous) & 8 \\
graph & $m \in \{3, 10\}$ & $9 \times 2$ & 18 \\
national & $n = 400$ & 9 & 9 \\
\midrule
total & & & 260 \\
\bottomrule
\end{tabular}
\end{center}
In the core every combination of $(n, T, \omega)$ exists for each of the eight scenario-variants with three regimes, and S2 enters once per $(n, T)$; no cell belongs to one scenario only. The reference cell belongs to the core, and the other blocks are counted without it. The first pass runs $M = 10$ replications, 2600 runs, to measure the actual cost of every cell before the final number of replications is fixed; its replications are the first ten of any later run, with the same seeds, and the files of the runs compose on the key (cell, replication). The ten additional scenario-variants are run at the reference levels only.

A full factorial of the five margins, with $n \in \{40, 100, 200, 400\}$, $m \in \{3, 5, 10\}$ and both balances, would have 1764 cells; it was started and stopped after 83 of its 17\,640 runs, since at the measured mean of about 130 seconds a run it would have taken some 650 core-hours.

\paragraph{How the blocks are read.} The core maps the performance of the procedure along the three factors that govern the information available to it: $n$, how many locations each regime has; $T$, how much evidence accumulates for each label; $\omega$, whether space helps, since at $\omega = 0$ the regimes occupy the same area and can be told apart by their parameters only. The other three blocks are checks of robustness at a representative configuration, the reference cell. Each tests a hypothesis with an expected direction, and each contains the two hard scenario-variants, S1-weak-\texttt{shr} and S3-beta-\texttt{shr}, where an effect would show; in the others the reference cell is at the ceiling.
\begin{description}[leftmargin=1.5em]
  \item[Balance.] At $1:2:3$ and $n = 100$ the regimes have 17, 33 and 50 locations. The likelihood gain of a regime grows with the number of its locations, so the smallest regime carries half the evidence of a balanced one; and its border locations have more neighbours in the larger regimes, towards which the Potts term pulls them. The expected answer: nothing changes in the well separated scenario-variants, while in the hard ones the small regime may be absorbed and $k$ underselected. The block says whether the conclusions of the core hold when the regimes differ in size.
  \item[Graph.] The graph of the penalty is a choice of the analyst (Section~2.5(a) of the manuscript), and the manuscript argues that under the automatic scale the factor $c$ divides by the mean degree, so that raising $m$ spreads the same total pull over more and longer edges rather than increasing it. The block tests that claim. The expected answer: results nearly invariant to $m$; a sparse graph ($m = 3$) is more local and lets small islands of labels survive, a dense one ($m = 10$) smooths across the borders over longer edges. It also measures the cost, which grows with $m$.
  \item[National scale.] $n = 400$ is the size of a national monitoring network, and the scale at which the theory of the supplement (Section~S2) predicts a risk. At $T = 120$ and $q = 8$ parameters per regime, the gain a superfluous fourth regime can extract from the search over partitions is bounded by $2n\log(4/3)$, which exceeds the BIC increment $q\log(nT)$ from $n = 200$ on:
\begin{center}
\begin{tabular}{rrr}
\toprule
$n$ & bound on the spurious gain, 3 to 4 regimes & BIC increment for one more regime \\
\midrule
100 & 58 & 75 \\
200 & 115 & 81 \\
400 & 230 & 86 \\
\bottomrule
\end{tabular}
\end{center}
  The bound is conservative, since neighbouring partitions share most of their labels, so it need not bite. The block, together with $n = 200$ in the core, says whether the BIC overselects on wide networks in practice, and so whether the extended BIC with $\gamma = 1$ is needed as a check there. It also measures the cost at that scale.
\end{description}

\paragraph{The limit of the design.} Varying one margin at a time does not show interactions: it cannot say whether the imbalance weighs more at $n = 40$, where the smallest regime has seven locations, or whether the graph matters more at $\omega = 0$, where borders are everywhere. The results of the three blocks are therefore statements of robustness at the reference configuration, not in general, and the paper must present them as such. This is the usual compromise when the full factorial cannot be afforded \citep{MorrisWhiteCrowther2019}. Should the interactions be asked for, the three blocks can be repeated at the hardest cell of the core, $n = 40$, $T = 60$, $\omega = 0$, some 35 cells more.

\paragraph{What is fitted.} Every replication runs the procedure a user would run: the grid $k \in \{1,\dots,4\}$ and $\phi \in \{0, 0.025, 0.05, 0.1, 0.2, 0.5, 1\}$ through \texttt{SCSTEM\_Infocrit}, and the selection rule \texttt{SCSTEM\_Select} with its default band $[0.025, 0.2]$: $k$ by the modal BIC over the band, the pooled fit competing, then $\phi$ by the smallest BIC at that $k$ over the whole grid. Every penalized fit starts from the unpenalized fit at the same $k$, and the scale of its penalty is measured there (Section~4.1 of the manuscript and the supplement). Selection is read off the selected fit; recovery of the partition and of the parameters is read off the fit at the true number of regimes and $\phi = 0.05$, so that a poor estimate is not confounded with a poor selection; the pooled fit is the benchmark. All of them are members of the grid.

\paragraph{What is recorded.} One row per replication with the selected $(k, \phi)$, the Adjusted Rand Index and the share of correctly assigned locations at the selected and at the true $k$, the error of the fitted signal against the conditional mean for the fit at the true $k$ and for the pooled fit, and the time; one row per replication, regime and parameter with the truth beside the estimate after aligning the estimated regimes on the true ones; one row per replication and location; one row per replication and grid point with the criteria and the Adjusted Rand Index of its partition, so that the recovery at any penalty, or under another selection rule, can be read without refitting; and the full observations for the first replications.

\paragraph{Replications.} $M = 10$ per cell in the first pass, with seeds derived from the replication index, so that a cell is reproducible on its own and replication $r$ uses the same locations and the same covariate in every scenario-variant.

\paragraph{Cost.} Measured on one replication of each of the 19 cells at the reference margins, with the grid and the rule above, run on 13 cores: one replication of the whole design takes 14 minutes of core time, a mean of 43 seconds per cell against a median of 33, with a tail in the cells where the EM works hardest (83 and 118 seconds for S3-Seta in the two variants, 98 for S3-beta in the independent one). One hundred replications take about 23 core-hours, 1.8 hours of wall time on 13 cores of the machine where it was measured. A single replication is a rough guide for a cost this heavy-tailed; the script recalibrates the estimate on every replication recorded. Adding a level to a margin multiplies the cost by the number of levels. The cost grows with $n$ far less than the factor 64 that $n^3$ would give from $n = 100$ to $n = 400$: one replication measured at every $(n, T)$ costs about 2 to 4 times more for S1-strong-\texttt{ind}, and 3 to 12 times more for S2, whose grid is dominated by the fits on the whole network; from $T = 120$ to $T = 365$ it costs two to three times more. S3-Seta-\texttt{shr} has a heavy tail at large $n$ and $T$: 21 minutes for one replication at $n = 200$, $T = 365$, and 66 at $n = 400$, $T = 365$. Priced on these measurements, with $n = 40$ priced as $n = 50$, one replication of the 260 cells of the experiment takes about 5.8 core-hours, and the first pass of ten replications about 58 core-hours, some 8 hours on the 7 cores of the machine that runs it. Cells with $\omega = 0$ cost about twice as much as the others in the stopped full factorial, so the first pass can take 10 to 12 hours.

%%=============================================================================
\section{Development extensions (not for the paper)}\label{sec:dev}
%%=============================================================================

Two extensions came out of the design of S3 and of the shared variants. They are implemented as functions outside the package, in the development area of the repository, are not used by the current version of the package, and are not mentioned in the paper.

\subsection{Constrained estimation: blocks common across regimes}\label{sec:dev:constr}

\paragraph{The constraint.} Split $\Psi_g = (\psi_c, \psi^s_g)$ into a block $\psi_c$ common to all regimes and the specific rest. The log-likelihood keeps the form \eqref{eq:loglik} on a restricted parameter space,
\[
  \ell(\psi_c, \psi^s_1, \dots, \psi^s_k) = \sum_{g=1}^k \ell_g(\psi_c, \psi^s_g),
  \qquad R\, \mathrm{vec}(\Psi) = 0,
\]
with $R$ the $(k-1)d_c \times k d$ matrix of the differences $\psi_{c,g} - \psi_{c,1}$. The number of parameters falls from $k_{\mathrm{eff}}(d_c + d_s)$ to $d_c + k_{\mathrm{eff}} d_s$. Given the partition, the free and the constrained models are nested, and $2(\ell_U - \ell_R)$ is asymptotically $\chi^2_{(k-1)d_c}$; since the partition is estimated, the reference distribution is taken from the bootstrap, or the two models are compared by an information criterion.

\paragraph{The admissible blocks.} \Stem{} parametrizes the error covariance as $\sigma^2_\omega(b\,I + C_\theta)$ with $b = \sigma^2_\epsilon/\sigma^2_\omega$, and estimates $(\sigma^2_\omega, \log b, \log\theta)$. The five blocks are the ones whose equality is linear in that parametrization: beta (the whole vector), $G$, $\Sigma_\eta$, $\theta$, and error ($\sigma^2_\epsilon$, $\sigma^2_\omega$ and $\theta$ together, equivalent to common $\sigma^2_\omega$, $b$ and $\theta$). Blocks can be combined. A common $\sigma^2_\epsilon$ alone or a common $\sigma^2_\omega$ alone is excluded: it is the nonlinear restriction $\sigma^2_{\epsilon} = b_g \sigma^2_{\omega,g}$ for all $g$, which would require reparametrizing the Newton--Raphson step in $(\log\sigma^2_\epsilon, \log\sigma^2_\omega, \log\theta)$ with new derivatives.

\paragraph{The EM algorithm.} The E-step is unchanged: one Kalman smoother per regime, since the latent processes are independent, giving the regime sufficient statistics $S^{(g)}_{00}, S^{(g)}_{10}, S^{(g)}_{11}$, $\Omega_g = \sum_t \E[e_{t,I_g} e_{t,I_g}^\top]$ and the generalized least squares blocks. In the M-step the specific blocks are updated regime by regime as now, and the common ones pool the statistics of all regimes:
\begin{align*}
  \beta &= \Bigl[\sum_g \sum_t X_{g,t}^\top \Sigma_{e,g}^{-1} X_{g,t}\Bigr]^{-1} \sum_g \sum_t X_{g,t}^\top \Sigma_{e,g}^{-1}\bigl(\hat z_{g,t} - K_g m^{(g)}_t\bigr), \\
  G_{jj} &= \frac{\sum_g [S^{(g)}_{10}]_{jj} / [\Sigma_{\eta,g}]_{jj}}{\sum_g [S^{(g)}_{00}]_{jj} / [\Sigma_{\eta,g}]_{jj}}
     \quad\Bigl(\text{and } G = \textstyle\sum_g S^{(g)}_{10}\bigl(\sum_g S^{(g)}_{00}\bigr)^{-1} \text{ if } \Sigma_\eta \text{ is common too}\Bigr), \\
  \Sigma_\eta &= \frac{1}{kT}\sum_g \Bigl[S^{(g)}_{11} - G_g S^{(g)\top}_{10} - S^{(g)}_{10} G_g^\top + G_g S^{(g)}_{00} G_g^\top\Bigr],
\end{align*}
a common $\theta$ by Newton--Raphson on $\log\theta$ jointly with the specific $\log b_g$, with gradient and Hessian summed over the regimes, and a common error block by $\sigma^2_\omega = \sum_g \mathrm{tr}[(bI + C^{(g)}_\theta)^{-1}\Omega_g] / (T\sum_g n_g)$ and Newton--Raphson on $(\log b, \log\theta)$ with summed derivatives. When the update of a common block depends on a specific one ($G$ common with $\Sigma_{\eta,g}$ specific) the updates are conditional maximizations in sequence, an ECM algorithm, which keeps the monotonicity of the EM for a fixed partition; the alternation as a whole remains non-monotone, as it is now.

\paragraph{The alternating algorithm.} The parameter step, now $k$ independent fits, becomes one joint EM: at every iteration $k$ E-steps and one coupled M-step, with joint convergence. The assignment step keeps its form \eqref{eq:score} with $\beta_h \to \beta$ and $s^2_h \to s^2$ where these are common. The common blocks then drop out of the contrasts of the score: with $\beta$ and the error variances common, the score discriminates through $\hat y^{(h)}$ alone. With a common error variance the mechanism of degeneracy, the regime with the smallest variance attracting every location, also disappears from the score. The data-driven scale of the penalty adapts by itself, being computed on the spread of the scores. A ridge on a common $\beta$ has to be scaled on the total number of locations rather than on the size of a regime.

\paragraph{Where it lives.} External functions in the development area: the joint EM for a fixed partition, the alternating algorithm that uses it, and the count of parameters. The package is not modified; without common blocks the external functions must return the numbers of \texttt{SCSTEM\_Estimation}.

\subsection{The coupled model: dependence between regimes}\label{sec:dev:coupled}

\paragraph{The model.} Given the partition, stack the latent processes into one state $y_t = (y^{(1)}_t, \dots, y^{(k)}_t)^\top$:
\[
  z_t = X_t \beta + L\, y_t + e_t, \qquad y_t = \mathrm{diag}(G_1,\dots,G_k)\, y_{t-1} + \eta_t, \qquad \eta_t \sim N_k(0, \Sigma_\eta),
\]
with loading $L_{ig} = 1\{i \in I_g\}$, regression coefficients by regime through covariates interacted with the regime indicators, $\Sigma_\eta$ a full matrix whose off-diagonal entries are the covariances $\rho_{gh}\sigma_{\eta,g}\sigma_{\eta,h}$ of the shocks, and $e_t$ either one field over the network or fields by regime. The log-likelihood is that of one STEM model on all $n$ locations and no longer splits into regime terms.

\paragraph{Nesting.} SC-STEM is the case with $\Sigma_\eta$ diagonal and fields by regime, in which the likelihood factorizes into \eqref{eq:loglik}. The pooled STEM is the case with common parameters, $\Sigma_\eta$ of rank one with equal entries (one path) and one field. With the coupled model the comparison between pooled and SC-STEM becomes one between nested models, and the independence of the regimes becomes testable.

\paragraph{Cost, in three levels.}
\begin{enumerate}[leftmargin=2em]
  \item \emph{Error field fully common} ($\sigma^2_\epsilon$, $\sigma^2_\omega$, $\theta$ equal across regimes). Most of it exists in STEM: the Kalman filter on $n$ locations with a multi-dimensional state, the current error covariance, the regression by regime through interactions, and a full $\Sigma_\eta$ (\texttt{flag.Sigmaetadiag = FALSE}). One correction is needed: with $G$ diagonal and $\Sigma_\eta$ full, the update of $G$ in the current code is not the exact conditional maximizer, which solves $(\Sigma_\eta^{-1} \circ S_{00})\, g = \mathrm{diag}(\Sigma_\eta^{-1} S_{10})$ for $g = \mathrm{diag}(G)$. Then a wrapper for the alternating algorithm and a bootstrap generator with dependence between regimes.
  \item \emph{Error parameters by regime with one field.} The covariance becomes $\Sigma_e = D_\epsilon + S_\omega C_\theta S_\omega$ with $D_\epsilon = \mathrm{diag}(\sigma^2_{\epsilon,g_i})$ and $S_\omega = \mathrm{diag}(\sigma_{\omega,g_i})$, which needs a new covariance function and its derivatives: the same kind of work as the excluded nonlinear restriction.
  \item \emph{Inside the package}: a second mode of the parameter step, the criteria, the selection, the bootstrap and the printing to adapt, and a cost per iteration that grows from $\sum_g n_g^3$ to $n^3$, about nine times for three balanced regimes, for every partition the alternation visits.
\end{enumerate}
The agreed scope is level 1, as external functions in the development area.

\paragraph{A correction to the scope of level 1.} With the field restricted to be common, level 1 nests the pooled STEM but not SC-STEM, whose error covariance is block-diagonal with parameters by regime. Nesting SC-STEM as well requires the option of fields by regime; that option needs no new derivatives, because with a block-diagonal error covariance the M-step of the error separates by regime and reuses the existing one. Whether to include it in level 1 is to be decided.

\begin{thebibliography}{1}
\bibitem[Morris, White and Crowther(2019)]{MorrisWhiteCrowther2019}
Morris, T.~P., White, I.~R. and Crowther, M.~J. (2019).
\newblock Using simulation studies to evaluate statistical methods.
\newblock \emph{Statistics in Medicine}, 38(11), 2074--2102.
\newblock \href{https://doi.org/10.1002/sim.8086}{doi:10.1002/sim.8086}.
\end{thebibliography}

\end{document}
